\documentclass[11pt]{article}
\usepackage[margin=1in]{geometry}
\usepackage[T1]{fontenc}
\usepackage{lmodern}
\usepackage{microtype}
\usepackage{amsmath,amssymb}
\usepackage[hidelinks]{hyperref}

\newcommand{\F}{\mathcal F}
\newcommand{\G}{\mathcal G}
\newcommand{\N}{\mathbb N}
\newcommand{\C}{\mathbb C}

\title{A Literature-Implied Complex Subcase of Brech--Pi\~na's Dual-Isometry Question}
\author{Research note for run \texttt{fa\_banach\_001}}
\date{18 July 2026}

\begin{document}
\maketitle

\paragraph{Status.}
\textbf{Literature-implied answer (partial subcase).}
The countable complex-scalar version of Brech--Pi\~na's Question 6 has a
positive answer by Fakhoury's later theorem on isometries of
combinatorial-like Banach spaces.

\paragraph{Source question.}
Brech and Pi\~na define, for a compact hereditary family
$\F\subseteq[\kappa]^{<\omega}$ containing the singletons, the combinatorial
space $X_\F$ as the completion of $c_{00}(\kappa)$ for
\[
  \|x\|_\F=\sup_{s\in\F}\sum_{\alpha\in s}|x(\alpha)|.
\]
After proving their primal Banach--Stone theorem under the hypothesis
$[\kappa]^1\subseteq\overline{\F^{MAX}}\cap\overline{\G^{MAX}}$, they ask:
\[
\text{Is every isometry }T:X_\F^*\to X_\G^*
\text{ weak}^*\text{-weak}^*\text{ continuous?}
\]
Equivalently, is every such $T$ the adjoint of an isometry
$S:X_\G\to X_\F$?  This is Question 6 in Section 2 of
\cite{BP}.

\paragraph{Supporting theorem.}
Fakhoury works with countable combinatorial families on $\N$ and with the
more general spaces $X_{\F,\mathbf E}$.  In Corollary 2.15, specialized to
the complex scalar field and $\mathbf E=\ell_1$, she obtains that if
$\F_1,\F_2$ are compact combinatorial families, then every isometry
\[
  J:X_{\F_1}^*\longrightarrow X_{\F_2}^*
\]
is an $(\F_1,\F_2)$-compatible signed permutation of the coordinate
functionals.  This is also stated in Proposition 1.3(b) for one family:
if $\mathbb K=\mathbb C$ and $\F$ is compact, every isometry of $X_\F^*$ is
an $\F$-compatible signed permutation.  See \cite[Proposition 1.3 and
Corollary 2.15]{Fakhoury}.

\paragraph{Implication.}
A compatible signed permutation on the dual basis has the form
\[
  J e_i^*=\omega_i e_{\pi(i)}^*,\qquad |\omega_i|=1,\quad
  \pi(\F_1)=\F_2.
\]
The corresponding operator on the primal spaces sends
$e_{\pi(i)}$ to $\omega_i e_i$ (with the evident inverse convention) and is
an isometry because $\pi$ transports the family.  Therefore $J$ is the
adjoint of a primal isometry.  In particular it is
weak$^*$-weak$^*$ continuous.  Thus Brech--Pi\~na's Question 6 is answered
positively when $\kappa=\N$ and the scalar field is $\C$, even without the
extra assumption that every singleton lies in the closure of the maximal
sets.

\paragraph{Scope and limitations.}
This packet does not settle the real scalar case.  Fakhoury explicitly notes
that the analogous arbitrary-family real statement is false, citing work in
progress of Brech--Tcaciuc, but that warning does not by itself decide the
real case under Brech--Pi\~na's stronger closure-of-maximals hypothesis.  The
packet also does not address uncountable $\kappa$.

\paragraph{Search evidence.}
The target was checked against the run indexes for \texttt{2007.11071},
\texttt{Banach-Stone}, \texttt{combinatorial Banach}, \texttt{weak*}, and
\texttt{dual isometry}.  No prior packet for this exact paper was found.
Web and local-source searches for the source title, the displayed
weak$^*$-continuity question, and ``$X_\F^*$ combinatorial Banach isometry''
led to Fakhoury's arXiv:2305.13125.  Later papers on dual quasi-norms and
the Zoo of combinatorial Banach spaces did not appear to answer the
weak$^*$-continuity question.

\begin{thebibliography}{9}
\bibitem{BP}
C. Brech and C. Pi\~na,
\emph{Banach-Stone-like results for combinatorial Banach spaces},
Ann. Pure Appl. Logic 172 (2021), no. 8, Paper No. 102989;
arXiv:2007.11071.

\bibitem{Fakhoury}
M. Fakhoury,
\emph{Isometries of $p$-convexified combinatorial Banach spaces},
Studia Mathematica (2024);
arXiv:2305.13125.
\end{thebibliography}

\end{document}
